\section{Feature Engineering}

\subsection{Dataset Construction}

The dataset was constructed from 612 posts shared across six user categories. All possible category pairs were generated for each post, resulting in:

\[
\binom{6}{2}=15
\]

pairwise interactions per post. Each observation represents a user-category pair and its corresponding divergence score, producing a final dataset of:

\[
612 \times 15 = 9,180
\]

observations.

\subsection{Extracted Features}

The divergence score was used as the target variable. The predictive features were grouped into four categories:

\begin{itemize}
    \item \textbf{Engagement:} \textit{post\_likes}, \textit{post\_comments}, \textit{log\_likes}, \textit{log\_comments}
    \item \textbf{Temporal:} \textit{months\_since\_post}
    \item \textbf{User attributes:} \textit{verified\_encoded}, \textit{gender\_match}
    \item \textbf{Stance interactions:} eight one-hot encoded variables representing user stance-pair combinations
\end{itemize}